\documentclass[book.tex]{subfiles}
\begin{document}

\chapter{Hopfield Networks and Boltzmann Machines}
\label{chap:hopfield}
\noindent\rule{11cm}{0.4pt} \\
\textbf{Prerequisites:}  \autoref{chap:molecular_hamiltonian},   \\
\textbf{Difficulty Level:} ** \\
\noindent\rule{11cm}{0.4pt}
\vspace{5mm}

The Hopfield model and Boltzmann machine are two families of models inspired by neuroscience and statistical mechanics. These two modes are in fact mathematically equivalent and are tied to spin-glass models from statistical mechanics. \cite{leonelli2021effective}, \cite{agliari2020generalized}, \cite{agliari2020neural}, \cite{agliari2020tolerance}, \cite{fachechi2019dreaming}

\section{Hopfield Network}

The Hopfield network has cost function given by
\begin{align}
H(\sigma| \xi) &= -\frac{1}{2N} \sum_{\mu=1}^P \sum_{i,j=1}^{N,N}\xi^\mu_i \xi^\mu_j \sigma_i \sigma_j
\end{align}
Roughly, the Hopfield loss sums over random Ising masks of the model.

% " In this talk I am focusing on two paradigmatic models for storing and learning, that is, respectively, the Hopfield model and the Boltzmann machine. First, I will show their formal equivalence and then I will leverage this equivalence to improve our understanding and our mathematical control on machine learning. "

%References 

\end{document}